\documentclass[10pt,a4paper]{amsart}
\usepackage[margin=19mm]{geometry}
\usepackage{amsmath,amssymb,mathtools}
\usepackage[scr=rsfs,cal=euler]{mathalfa}
\usepackage{xfrac}
\usepackage[unicode,hidelinks,bookmarks=false]{hyperref}
\newcommand{\ie}{i.e.\ }
\newcommand{\cf}{cf.\ }
\newcommand{\resp}{respectively\ }
\newcommand{\Subsection}{\subsection}
\providecommand{\nobreakdash}{\nobreak\hbox{-}}
\setlength{\emergencystretch}{2em}
\multlinegap0pt
\usepackage{enumerate}
\usepackage[normalem]{ulem}
\usepackage{amssymb}
\newtheorem{thm}{Theorem}
\newtheorem{lem}{Lemma}
\newcommand{\R}{\mathbb{R}}
\def\az{\alpha}
\def\bint{{\ifinner\rlap{\bf\kern.35em--}
\int\else\rlap{\bf\kern.45em--}\int\fi}\ignorespaces}
\def\dist{{\mathop\mathrm{\,dist\,}}}
\def\ez{\epsilon}
\def\loc{{\mathop\mathrm{\,loc\,}}}
\title{Serrin's overdetermined theorem within Lipschitz domains}
\author{Hongjie Dong, Yi Ru-Ya Zhang}
\date{Source: arXiv:2509.05155v4 (CC BY 4.0)}
\begin{document}
\maketitle

\noindent\emph{Source and licence.} Serrin's overdetermined theorem within Lipschitz domains. Version arXiv:2509.05155v4 (CC BY 4.0), distributed by arXiv under the Creative Commons Attribution 4.0 International licence, http://creativecommons.org/licenses/by/4.0/, as recorded in the arXiv licence metadata for this exact version and shown on the abstract page https://arxiv.org/abs/2509.05155v4. Full licence notice: \texttt{licenses/arxiv-2509.05155-CC-BY-4.0.txt}.\par\smallskip

\noindent\textit{Editorial scope.} Counted unit: the article's lem DKP (source0.tex lines 559--568). The article's complete original proof of it is delivered with it (source0.tex lines 569--609, one \texttt{proof} environment of 41 lines). No mathematics is written, repaired or re-derived: the statement and the proof are the article's own lines, in the article's own notation, and no retained line is substituted or re-flowed.  Retained uncounted as the source-local context the proof needs: lines 308--310; lines 325--334.  Nothing was omitted from any retained span, so the package carries no omission record.  No retained line contains a citation, so the wrapper carries no bibliography and no bibliography entry.  No retained line contains a figure environment, an image-inclusion command or drawing code, so no image is delivered and the package declares no asset dependency.  Editorial formatting only, in addition: an AMS \texttt{amsart} preamble that restates the article's own declarations and macros used by the retained lines, namely the environments \texttt{thm}, \texttt{lem} on separate counters, together with the standard packages those lines need.  The retained environments print in this document as Theorem 1, Lemma 1 in order of appearance, whereas the article prints the same items with the numbering of its own full document.  The complete native source of the article as distributed in the arXiv e-print for this exact version is bundled unchanged as \texttt{sources/184539/source0.tex}.
\begin{equation}\label{weak wulff u}
   u \in W^{1,2}(\mathbb R^n), \quad  u=0\  \text{ a.e. in }\R^n\setminus \Omega,\quad \Delta_H u=\mathbf{c}H(-\nu)\mathscr{H}^{n-1}|_{\partial^*\Omega} - \mathbf{1}_{\Omega}\,dx. 
\end{equation} 

\begin{thm}\label{main thm 2}
    Let $\Omega$ be a Lipschitz domain with a local Lipschitz constant $L>0$, $K$ be a (bounded) uniformly convex set of $C^{3}$ class with the principle curvatures of $\partial K$ bounded from above by $\kappa$ and from below by $\kappa^{-1}>0$, 
    and $H$ be the associated Wulff potential.
    Assume that $u\in W^{1,\,2}(\mathbb R^n) $
    satisfies the anisotropic overdetermined system \eqref{weak wulff u}. Moreover, we suppose that $D^2u\in L^n$ and $|Du|\ge c_0>0$ in some neighborhood of $\partial \Omega$. Then there exists a constant $\ez_0=\ez_0(n,\,\kappa,\,\|K\|_{C^{3}})>0$ satisfying that, whenever
    $ L \le \ez_0,$
    one has
    $$u(x)=\frac{r^2-H_*^2(x)}{2n} \quad \text{ in  }  \Omega  \text{ for some $r>0$,}$$
where $H_*$ is the dual function of $H$. Moreover, this implies that $\Omega$ and $K$ are homothetic. 
\end{thm}

\begin{lem}\label{A DKP}
    Let $\Omega$ and $u$ be  as in Theorem~\ref{main thm 2}. Then 
    $$\mathcal A(x)= D^2V( Du(x))$$
    satisfies the $C\omega_0$-DKP condition in $\Omega$, i.e.
    $$\mu_\mathcal A(B_r(y)\cap \Omega)\le C \| D^2 u\|^{2}_{L^n(\mathcal N_\rho(\partial \Omega)\cap \Omega)}r^{n-1},\quad 0<r<\rho$$
    where, for some $\rho>0$, 
$$\mathcal N_\rho(\partial \Omega): =\{z\in \mathbb R^n\colon \dist(z,\,\partial\Omega)<\rho\}, \quad \omega_0=\|D^2u\|^2_{L^n(\mathcal N_\rho(\partial \Omega)\cap \Omega)},$$
    and
    $$0<C=C(n,\,c_0,\,L,\,\kappa,\,   \|D^3V\|_{L^\infty(\mathbb S^{n-1})}).$$
\end{lem}

\begin{proof}
    Let $y\in \partial\Omega$ and $x\in B_{r}(y)$.
    By differentiating the equation in \eqref{weak wulff u} in the $e_k$-direction, $1\le k\le n$, and letting $v=\partial_k u$, we obtain that $v$ satisfies, in the sense of distribution, 
    $${\rm div} (D^2V(Du) Dv)={\rm div}([H(Du)D^2H(Du)+DH(Du)\otimes DH(Du)]Dv )=0.$$
    According to the De Giorgi--Moser--Nash theory, this implies that $v\in C_\loc^{\alpha}$ with
    $$\alpha=\az(n,\,\kappa,\,\|K\|_{C^1})>0,$$
    where $\kappa$ is the curvature bound of $\partial K$ in Theorem~\ref{main thm 2}. 
    Then by summing over all $k$'s, as $|Du|\neq 0$ in a neighborhood of $\partial \Omega,$ it follows from the Schauder estimates in divergence form together with Poincar\'e inequality that, for any  $0<t\le \frac{d(x)} 4$ with $B_t(x)$ disjoint from $\Omega_0\subset\subset\Omega$, 
    \begin{align}
        \|D^2u\|_{L^\infty(B_t(x))}\le &\ C(n,\,\az,\,\|D^2V\|_{C^\az(\mathbb S^{n-1})})t^{-1} \inf_{{\bf a}\in\mathbb R^n}\bint_{B_{2t}(x)}|Du-{\bf a}|\,d\mathcal L^n\nonumber\\
        \le & \ C(n,\,\az,\,\|D^2V\|_{C^\az(\mathbb S^{n-1})}) \left(\bint_{B_{2t}(x)} |D^2 u|^n\, dx\right)^{\frac 2 n}. \label{bounded hessian}
    \end{align} 

    A direct computation yields
    $$D\mathcal A(z)= D^3V(Du(z))D^2u(z), $$
    and the $(-1)$-homogeneity of $D^3V$ further gives that,  for any $z\in B_{ {d(x)} /4}(x)$, 
    $$|D\mathcal A(z)|= \left|D^3V\left(\frac{Du}{|Du|}(z)\right)\frac{D^2u}{|Du|}(z)\right|\le c_0^{-1}\|D^3V\|_{L^\infty(\mathbb S^{n-1})}\|D^2u\|_{L^\infty(B_{ {d(x)} /4}(x))}.$$
    Therefore, by the Bescovitch covering theorem, we can find a collection of countably many balls 
    $$\{B_{d(x_i)/4}(x_i)\}_{x_i\in B_r(y)\cap \Omega},$$
    with uniformly bounded overlaps up to $C(n)$, covering $B_r(y)\cap \Omega$ so that,  via applying \eqref{bounded hessian} to $u$ on each $B_{d(x_i)/4}(x_i)$, 
\begin{align}
    & \mu_\mathcal A({B_r(y)\cap \Omega})\nonumber\\
    \le& \ C(n)\sum_{i} \mu_\mathcal A(B_{d(x_i)/4}(x_i))\nonumber\\
    \le & \ C(n)\sum_{i} |B_{d(x_i)/4}(x_i)| d(x_i) \|D^3V\|^2_{L^\infty(\mathbb S^{n-1})}\|D^2u\|^2_{L^\infty(B_{d(x_i)/4}(x_i))}\nonumber\\
    \le & \ C(n,\,c_0,\,\az,\, \|D^3V\|_{L^\infty(\mathbb S^{n-1})}) \sum_i  d(x_i)^{n-1} \left(\int_{B_{d(x_i)/2}(x_i)} |D^2 u|^n\, dx\right)^{\frac 2 n}.\label{mu estimate}
\end{align}
Then when $n=2$, via $d(x_i)\le r$ we directly get 
$$ \mu_\mathcal A({B_r(y)\cap \Omega})\le  C(c_0,\,\az,\,\|D^2V\|_{C^\az(\mathbb S^{1})},\,\|D^3V\|_{L^\infty(\mathbb S^{1})}) \| D^2 u\|^{2}_{L^2(\mathcal N_\rho(\partial \Omega)\cap \Omega)}r$$
as desired. 

When $n\ge 3$, we continue from \eqref{mu estimate} via H\"older's inequality that, for $\beta=\frac{n(n-1)}{n-2},$
\begin{align*}
    & \mu_\mathcal A({B_r(y)\cap \Omega})\\
    \le & \  C(n,\,c_0,\,L,\,\az,\, \|D^3V\|_{L^\infty(\mathbb S^{n-1})}) \left( \sum_i \int_{B_{d(x_i)/2}(x_i)} |D^2 u|^n\, dx\right)^{\frac 2 n}\left(\sum_id(x_i)^\beta\right)^{1-\frac 2 n}   \\
    \le & \  C(n,\,c_0,\,L,\,\az,\, \|D^3V\|_{L^\infty(\mathbb S^{n-1})}) \| D^2 u\|^{2}_{L^n(\mathcal N_\rho(\partial \Omega)\cap \Omega)}\left(\sum_id(x_i)^\beta\right)^{1-\frac 2 n}.
\end{align*}
Since there are at most $C(n,\,L)2^{l(n-1)} $-many balls in the collection with radius between $2^{-l} r$ and $2^{-l+1}r$ for each $l\in \mathbb N$, then
$$\left(\sum_id(x_i)^\beta\right)^{ 1-\frac 2 n}\le \left(\sum_l 2^{l(n-1)} 2^{-l(n-1)\frac{n}{n-2}}\right)^{ 1-\frac 2 n} r^{n-1}\le C(n)r^{n-1},$$
where  $-\frac{n}{n-2}+1<0$ was used. 
Thus, we obtain the desired estimate for $y\in \partial \Omega. $   
\end{proof}

\end{document}
